\documentclass[11pt,a4paper]{article}

\usepackage[utf8]{inputenc}
\usepackage[T1]{fontenc}
\usepackage{amsmath, amssymb}
\usepackage{graphicx}
\usepackage[margin=1in]{geometry}
\usepackage{hyperref}

\title{{name}}
\author{Your Name\\ \small Your Institution}
\date{\today}

\begin{document}

\maketitle

\begin{abstract}
    A short summary of the problem, what you did about it, and what you
    found. Aim for a paragraph a reader can understand on its own.
\end{abstract}

\section{Introduction}
\label{sec:introduction}

Set out the problem and why it matters, then say what this paper
contributes. Cite related work as you go~\cite{knuth1984texbook}.

\section{Method}
\label{sec:method}

Describe the approach in enough detail to be reproduced. Displayed
equations are numbered and can be referenced, like
Equation~\ref{eq:example}:

\begin{equation}
    \label{eq:example}
    E = mc^2
\end{equation}

\section{Results}
\label{sec:results}

Present the findings. Refer to figures and tables by label rather than
position, so they stay correct when the layout changes.

\section{Conclusion}
\label{sec:conclusion}

Restate what was learned and name the questions left open.

\bibliographystyle{plain}
\bibliography{references}

\end{document}
